\section{算法与编程}

\answer{AG4-1}
A 组有 $n_A=3$、$\mu_A=4$、$M_A=9+0+9=18$，故 B 组有 $n_B=2$，总和为 $5\cdot4-3\cdot4=8$，均值4。两组均值相同，组间平方和为零，所以 $M_B=26-18=8$。合并总体方差为 $26/5$，样本方差为 $26/4=13/2$。

B 组的两个数可写成 $4-u,4+u$，中心平方和为 $2u^2=8$，故 $|u|=2$。因此无序数对唯一为 $\{2,6\}$，但原先的先后顺序不能由摘要确定。通常均值和平方和不能恢复整组数据；这里的唯一性依赖 B 组恰有两个数。

\answer{AG4-2}
区间 $[0,3]$ 的值为 $(4,2,3,5)$，极差3，长度4。区间 $[3,6]$ 也有长度4且极差3，所以按最早起点选择 $[0,3]$。

仅有的三个长度5区间的极差依次为6、6、5，都超过3。任何长度超过5的区间都包含一个长度5区间，极差不会更小，所以不存在更长可行解。双单调队列扫描得到相同结果；这个排除论证独立于程序状态，直接证明长度4最优。

\answer{AG4-3}
期限为2时，各终点的最好合法候选如下：终点1为 $1-80/120=1/3$；终点2为 $1-100/120=1/6$；终点3只能使用80、100两个峰值，最好为 $1-70/100=3/10$；终点4最好为 $1-90/100=1/10$。所以最大值为 $1/3$，区间 $[0,1]$。

不设期限时，120到70产生 $1-70/120=5/12$，区间 $[0,3]$，且它是全局最优。这个区间持续3个时间单位，超出前一问题的可行域。保留过期峰值会把后一个答案误用于前一个问题。

\answer{AG4-4}
从80到60恰好下降 $1/4$，对应区间 $[2,4]$，持续2。相邻四次降幅为 $6/100$、$14/94$、$1/80$、$19/79$，都严格小于 $1/4$，所以不存在持续1的可行解，最短持续时间确为2。

对固定起点 $i$，扫描至第一个达标终点 $j$ 后，它在未来只能形成更长区间。若它本来已被一个更晚且价格不低的起点支配，删除也不会损失更好的结果。因此永久删除有明确的目标依据；这种删除规则不能不加证明地用于最大回撤比例问题。

\par\Needspace{8\baselineskip}
\subsection*{小节练习}
\answer{OA31}
用长度为 $N$ 的循环缓冲区保存窗口，以便在满窗时取出最旧值。设当前窗口大小为 $k=\min(t,N)$，维护
\[
 S=\sum_{\text{window}}x_i,\qquad Q=\sum_{\text{window}}x_i^2.
\]
未满窗时，新值 $b$ 到来后令 $S'=S+b$、$Q'=Q+b^2$、$k'=k+1$。满窗时设移出的值为 $a$，则
\[
 S'=S-a+b,\qquad Q'=Q-a^2+b^2,
 \qquad k'=N.
\]
输出
\[
 \mu'=\frac{S'}{k'},\qquad
 v'=\frac{Q'}{k'}-(\mu')^2.
\]
恒等式 $\sum(x_i-\mu)^2=Q-k\mu^2$ 证明了输出正确。每次只需常数次算术操作，故在单位代价算术模型下单次更新时间为 $O(1)$；缓冲区占 $O(N)$ 空间。$N=1$ 时方差恒为零。

当数值很大而波动很小时，$Q/k-\mu^2$ 会有浮点消去误差。可以改维护中心平方和 $M=\sum(x_i-\mu)^2$。满窗的更新为
\[
 \begin{aligned}
 \mu'&=\mu+\frac{b-a}{N},\\
 M'&=M+(b-a)(b+a-\mu-\mu'),\qquad v'=M'/N.
 \end{aligned}
\]
第二式由 $M=Q-N\mu^2$ 和 $N(\mu'-\mu)=b-a$ 展开即得。未满窗可用增量公式：令 $\delta=b-\mu$，则
\[
 \mu'=\mu+\frac{\delta}{k+1},\qquad
 M'=M+\delta(b-\mu').
\]
空窗的第一个值直接初始化为 $\mu=b,M=0$。中心更新减少了反复相减两个大二阶矩的风险，但浮点误差仍可能累积，必要时可从缓冲区重新计算。若需要分母为 $k-1$ 的样本方差，在 $k>1$ 时输出 $M/(k-1)$；是否无偏还取决于抽样模型。

\answer{OA32}
维护一个双端队列，保存候选下标。队列内下标严格递增，对应数值严格递减。新值 $x_t$ 到来时依次执行：
\begin{enumerate}
\item 从队首删除所有小于 $t-N+1$ 的下标，即已离开窗口的元素。
\item 只要队列非空且队尾的值不大于 $x_t$，就删除队尾。
\item 将下标 $t$ 加入队尾，输出队首对应的值。
\end{enumerate}
删除队尾是安全的：若旧下标 $i<t$ 满足 $x_i\le x_t$，在任何同时包含它们的后续窗口内，$i$ 都不优于 $t$；而 $t$ 更晚过期。因此 $i$ 不会再成为必须保留的最大值候选。过期元素也不能影响当前窗口。更新后值按递减排列，队首就是当前最大值。相等值时保留较新的下标，因其有效期更长；若另需最早最大值下标，可把第二步的“不大于”改成“严格小于”。

每个下标只入队一次，至多出队一次。因此处理 $T$ 个值的总操作数为 $O(T)$，每次更新的摊还时间为 $O(1)$，队列空间为 $O(N)$。但一次很大的新值可以删除队列中几乎全部元素，所以单次最坏时间是 $O(N)$，不能把摊还界误写为逐次最坏界。$N=1$ 时输出就是新值本身。

\answer{OA33}
初始化历史峰值 $H=p_1$、峰值下标 $i_H=1$、答案 $D=0$。对 $j=2,\ldots,n$ 依次执行：先计算 $d=1-p_j/H$；若 $d>D$，更新 $D=d$ 并保存区间 $(i_H,j)$；然后若 $p_j>H$，更新 $H=p_j$、$i_H=j$。若始终没有正跌幅，就返回 $D=0$，且无需返回正回撤区间。

处理 $j$ 之前，$H=\max_{i<j}p_i$。对固定的正数 $p_j$，函数 $1-p_j/p_i$ 随正数 $p_i$ 增大而增大，所以以 $j$ 为终点的最大跌幅恰为 $1-p_j/H$。算法对每个终点都检查了它的最优起点，再与零取最大值，因此得到全局最大回撤。先计算再更新峰值保证了起点严格早于终点。

每个价格只处理一次，时间为 $O(n)$；除了几个数值和下标外不需要额外存储，辅助空间为 $O(1)$。使用严格的 $d>D$ 更新会保留最早出现的最大回撤终点，峰值相等时不更新则保留最早峰值下标。

回撤的分母必须是此前峰值。例如价格从 $100$ 跌到 $60$，回撤为 $40\%$。表达式 $H/p_j-1$ 算的是从低点恢复到峰值所需的涨幅，在本例中为 $2/3$，不是回撤。若用累计峰值 $H_j=\max_{i\le j}p_i$ 表示，同一答案可写为
\[
 \boxed{D=\max_{1\le j\le n}\left(1-\frac{p_j}{H_j}\right)}.
\]

